% 06 — Reparse / Junction：从零到不展开 + 可选重建

\section{起点：把 junction 当普通目录递归}

Windows **directory junction**（\winapi{IO\_REPARSE\_TAG\_MOUNT\_POINT}）在资源管理器中「看起来像文件夹」，实际是 reparse point。
若 walk 跟随进入 \texttt{D:/}，一次备份 \texttt{C:/project/link} 可能扫进整盘外置硬盘——\textbf{灾难性} scope 膨胀。

\begin{evolbox}[GAP-WIN-03 双轨策略]
\textbf{扫描层}：发现 reparse 目录 → 记 \texttt{reparse\_junction} issue，\textbf{不递归}，但采集 \texttt{reparse\_tag/target}。\\
\textbf{VSS 层}：\texttt{ProbeJunctionVolumes} 有限深度探测 target 卷，纳入 snapshot set（专册 09）。\\
\textbf{恢复层}：\texttt{ReparseRestorePolicy} — \texttt{kSkip} 或 \texttt{kRecreate}。
\end{evolbox}

\section{Reparse Point 类型}

ebbackup 重点支持 \winapi{IO\_REPARSE\_TAG\_MOUNT\_POINT}（目录联接）。
符号链接（\winapi{IO\_REPARSE\_TAG\_SYMLINK}）在 scan 层按 symlink 类型处理；junction 是\textbf{目录} + reparse 的组合。

持久化字段：
\begin{itemize}[nosep]
  \item \texttt{reparse\_tag} → \manifestflag{kMetaReparse}；
  \item \texttt{reparse\_target}（UTF-8，规范化后）→ \manifestflag{kMetaReparseTarget}。
\end{itemize}

\section{采集：FSCTL\_GET\_REPARSE\_POINT}

\texttt{CaptureReparseTarget}（\filepath{win\_meta.cc}）在 \texttt{CaptureWinMetaFromPath} 内调用：

\begin{lstlisting}[caption={CaptureReparseTarget (win\_meta.cc)}]
DeviceIoControl(h, FSCTL_GET_REPARSE_POINT, nullptr, 0,
                buf.data(), buf.size(), &ret, nullptr);

const auto* rdb = reinterpret_cast<const EbReparseDataBuffer*>(buf.data());
if (rdb->ReparseTag == IO_REPARSE_TAG_MOUNT_POINT) {
  const WCHAR* sub_name = ...;  // SubstituteName from PathBuffer
  std::wstring substitute(sub_name, sub_chars);
  *out_utf8 = WideToUtf8(NormalizeSubstituteName(std::move(substitute)));
}
\end{lstlisting}

\subsection{Substitute 名规范化}

Windows 存储的 substitute 常带 \texttt{\textbackslash??\textbackslash} 或 \texttt{\textbackslash DosDevices\textbackslash} 前缀：

\begin{lstlisting}[caption={NormalizeSubstituteName (win\_meta.cc)}]
if (wpath.rfind(L"\\??\\", 0) == 0) wpath = wpath.substr(4);
else if (wpath.rfind(L"\\DosDevices\\", 0) == 0) wpath = wpath.substr(12);
\end{lstlisting}

manifest 存\textbf{规范化后} UTF-8，恢复时再结合 restore 根目录重算绝对路径。

\section{扫描策略：不跟随联接}

\begin{lstlisting}[caption={Walk 层 junction 截断 (scan\_entry.cc)}]
if (IsReparseDirectory(entry.path())) {
  RecordIssue(added.absolute_path, "reparse_junction", &out->issues);
  winmeta::CaptureWinMetaFromPath(added.absolute_path, &added);
  continue;  // 不 push 子目录到 walk stack
}
\end{lstlisting}

\begin{figure}[htbp]
\centering
\begin{tikzpicture}[node distance=0.9cm]
  \node[wbox] (scan) {ScanDirectory DFS};
  \node[wdata, below left=of scan] (junction) {junction 目录};
  \node[wdata, below right=of scan] (normal) {普通子目录};
  \node[wbox, below=1.2cm of junction] (meta) {tag + target → manifest};
  \node[wdata, below=1.2cm of normal] (recurse) {继续递归};
  \draw[warr] (scan) -- (junction);
  \draw[warr] (scan) -- (normal);
  \draw[warr] (junction) -- (meta);
  \draw[warr] (normal) -- (recurse);
\end{tikzpicture}
\caption{Junction：采集元数据但不 recurse}
\end{figure}

避免的问题：
\begin{itemize}[nosep]
  \item 意外备份整盘；
  \item 与 VSS 多卷闭包逻辑冲突；
  \item 循环 junction 导致 scan 深度爆炸。
\end{itemize}

\section{恢复：ReparseRestorePolicy}

\begin{longtable}{@{}llp{7cm}@{}}
\toprule
模式 & CLI & 行为 \\
\midrule
\texttt{kSkip} & 默认 & 创建普通空目录或跳过 reparse 属性 \\
\texttt{kRecreate} & \texttt{reparse=recreate} & \texttt{RecreateReparsePoint} \\
\bottomrule
\end{longtable}

\texttt{restore\_engine.cc} 对 \texttt{file\_type==kDirectory \&\& reparse\_tag!=0} 优先分支：

\begin{lstlisting}[caption={RecreateReparsePoint 双路径 (win\_meta.cc)}]
// 1) 优先 cmd mklink /J（用户熟悉、权限清晰）
std::system("cmd /c mklink /J \"link\" \"target\"");

// 2) fallback: 手工 REPARSE_DATA_BUFFER + FSCTL_SET_REPARSE_POINT
DeviceIoControl(h, FSCTL_SET_REPARSE_POINT, &rdb, ...);
\end{lstlisting}

\section{目标路径重映射}

恢复时 \texttt{reparse\_target} 可能是：

\begin{itemize}[nosep]
  \item \textbf{绝对}：\texttt{D:/data} — 结合 restore 映射规则替换盘符；
  \item \textbf{相对}：相对 junction 父目录解析为 \texttt{GetFullPathNameW}。
\end{itemize}

开发机 \texttt{C:} → 恢复机 \texttt{D:} 场景：避免 hardcode 源机器盘符。

\section{VSS 多卷闭包（交叉引用）}

\texttt{vss\_volume\_closure.cc} 的 \texttt{ProbeJunctionVolumes}：

\begin{lstlisting}[caption={ProbeJunctionVolumes（概念）}]
// BFS 源树，深度 junction_probe_depth（默认 2）
// IO_REPARSE_TAG_MOUNT_POINT → 解析 target → ResolveVolumeForPath
// 跨卷 volume_key 加入 ComputeVolumeClosure 结果
\end{lstlisting}

\textbf{扫描不展开} $\neq$ \textbf{VSS 不包含目标卷}：快照可含 D:，但 manifest 只含用户指定的 C: 树 + junction 元数据。

\section{REPARSE\_DATA\_BUFFER 结构}

\begin{figure}[htbp]
\centering
\begin{tikzpicture}[
  cell/.style={rectangle, draw, minimum width=2.2cm, minimum height=0.45cm, font=\tiny, align=center}
]
  \node[cell, fill=blue!10] (tag) {ReparseTag};
  \node[cell, fill=green!10, right=0.1cm of tag] (len) {DataLength};
  \node[cell, fill=yellow!10, right=0.1cm of len] (sub) {SubstituteName};
  \node[cell, fill=orange!10, right=0.1cm of sub] (print) {PrintName};
  \node[below=0.4cm of tag, font=\scriptsize, text width=10cm, align=center] {
    ebbackup 持久化 Substitute 解析结果；PrintName 仅用于 UI，恢复以 Substitute 为准
  };
\end{tikzpicture}
\caption{Junction reparse data 概念字段}
\end{figure}

自研 \texttt{EbReparseDataBuffer} 布局与 WDK 文档对齐，避免 SDK 版本差异。

\section{Restore 排序}

\texttt{FileTypeOrder} 保证\textbf{目录}（含 reparse 目录）在文件之前创建——\texttt{RecreateReparsePoint} 需要父目录存在。

\section{测试}

\begin{itemize}[nosep]
  \item \filepath{win\_reparse\_scan\_test.cc} — scan 不进入 target；
  \item \filepath{win\_meta\_test.cc} — tag/target roundtrip；
  \item \filepath{vss\_volume\_closure\_test.cc} — junction 卷发现。
\end{itemize}

\section{明确不做的 reparse 类型}

\begin{warnbox}[非 MOUNT\_POINT]
OneDrive/cloud filter（\winapi{IO\_REPARSE\_TAG\_CLOUD} 等）、DFS 专用 tag：当前仅采集 tag，\texttt{kRecreate} 可能失败并记 issue。
\end{warnbox}
